%% fig-sequence.tex — Slide 5 «Як працює» ★ HAPAX LEGOMENON (the one dense slide).
%% SOURCE: adapted from chapter2.tex fig:auth_sequence (lines 223-305).
%% DO NOT MODIFY chapter2.tex — this is a standalone copy, restyled for projection.
%%
%% Trimming vs. thesis original:
%%   - «Реєстрацію» block (r1-r4) DROPPED — mentioned verbally only.
%%   - «Вхід (2FA)» (l1-l4) + «Встановлення тунелю» (v1-v4) KEPT.
%%   - Arrow labels shortened to ≤3 tokens (no JSON payloads).
%%   - Pastel actor fills → our actor style (clInk border, white fill).
%%   - arr/retarr → msg/msgret; double tunnel → tunnel style.
%%   - No callout on v2 — presenter covers verbally.
%%   - Only TTL annotations [5 хв] / [24 год] on slide (R2 rule).
%% Labels are folded INTO \draw commands (node[midway,above]{...}) — avoids
%% the origin-collapse bug that occurs with standalone \node[midway] at (x,y).
%% Assumes styles.tex pre-loaded (calc + arrows.meta needed in outer preamble).
\begin{tikzpicture}[x=1cm, y=1cm]
%%
%% ── Lifeline x-positions (matches thesis proportions) ────────────────────────
\def\xC{0}    % Клієнт
\def\xA{5.0}  % API Server
\def\xW{9.5}  % WireGuard
%%
%% ── Actor boxes — deck palette (no pastel fills) ─────────────────────────────
\node[actor, minimum width=2.6cm] (C) at (\xC, 0) {Клієнт};
\node[actor, minimum width=2.6cm] (A) at (\xA, 0) {API Server};
\node[actor, minimum width=2.6cm] (W) at (\xW, 0) {WireGuard};
%%
%% ── Lifelines (9.5 cm — covers all exchanges with room to spare) ─────────────
\draw[lifeline] (C.south) -- ++(0,-9.5cm);
\draw[lifeline] (A.south) -- ++(0,-9.5cm);
\draw[lifeline] (W.south) -- ++(0,-9.5cm);
%%
%% ── Block label: Вхід (2FA) ──────────────────────────────────────────────────
\node[font=\small\sffamily\bfseries, text=clGrey, align=right]
  at (-1.8, -0.55) {\textbf{Вхід (2FA)}};
%%
%% ── l1: login (пароль)  C → A ────────────────────────────────────────────────
\coordinate (l1s) at ($(C.south)+(0,-0.9cm)$);
\coordinate (l1e) at ($(A.south)+(0,-0.9cm)$);
\draw[msg] (l1s) -- (l1e)
  node[midway, above, font=\small\sffamily] {login (пароль)};
%%
%% ── l2: проміжний токен  A → C  +  TTL [5 хв] ────────────────────────────────
\coordinate (l2s) at ($(A.south)+(0,-1.9cm)$);
\coordinate (l2e) at ($(C.south)+(0,-1.9cm)$);
\draw[msgret] (l2s) -- (l2e)
  node[midway, above, font=\small\sffamily] {проміжний токен};
\node[font=\small\sffamily\bfseries, text=clAccent, right=0.2cm]
  at ($(A.south)+(0,-1.9cm)$) {\textbf{[5 хв]}};
%%
%% ── l3: TOTP  C → A ──────────────────────────────────────────────────────────
\coordinate (l3s) at ($(C.south)+(0,-3.1cm)$);
\coordinate (l3e) at ($(A.south)+(0,-3.1cm)$);
\draw[msg] (l3s) -- (l3e)
  node[midway, above, font=\small\sffamily] {TOTP};
%%
%% ── l4: повний токен  A → C  +  TTL [24 год] ─────────────────────────────────
\coordinate (l4s) at ($(A.south)+(0,-4.1cm)$);
\coordinate (l4e) at ($(C.south)+(0,-4.1cm)$);
\draw[msgret] (l4s) -- (l4e)
  node[midway, above, font=\small\sffamily] {повний токен};
\node[font=\small\sffamily\bfseries, text=clAccent, right=0.2cm]
  at ($(A.south)+(0,-4.1cm)$) {\textbf{[24 год]}};
%%
%% ── Block label: Тунель — moved down to align with «сесія» row (y=-5.35) ──────
%% Was at y=-4.95 (too close to «повний токен» in the 2FA phase).
\node[font=\small\sffamily\bfseries, text=clGrey, align=right]
  at (-1.8, -5.35) {\textbf{Тунель}};
%%
%% ── v1: сесія  C → A ─────────────────────────────────────────────────────────
\coordinate (v1s) at ($(C.south)+(0,-5.4cm)$);
\coordinate (v1e) at ($(A.south)+(0,-5.4cm)$);
\draw[msg] (v1s) -- (v1e)
  node[midway, above, font=\small\sffamily] {сесія};
%%
%% ── v2: wg set wg0 peer …  A → W  (accent, the KEY step) ────────────────────
\coordinate (v2s) at ($(A.south)+(0,-6.5cm)$);
\coordinate (v2e) at ($(W.south)+(0,-6.5cm)$);
\draw[msg, draw=clAccent, line width=1.5pt] (v2s) -- (v2e)
  node[midway, above, font=\small\ttfamily, text=clAccent]
    {\texttt{wg set wg0 peer \ldots}};
%%
%% ── v3: дані сесії  A → C ────────────────────────────────────────────────────
\coordinate (v3s) at ($(A.south)+(0,-7.6cm)$);
\coordinate (v3e) at ($(C.south)+(0,-7.6cm)$);
\draw[msgret] (v3s) -- (v3e)
  node[midway, above, font=\small\sffamily] {дані сесії};
%%
%% ── v4: тунель  C ↔ W  (two parallel arrows = established bidirectional tunnel)
\coordinate (v4s) at ($(C.south)+(0,-8.7cm)$);
\coordinate (v4e) at ($(W.south)+(0,-8.7cm)$);
\draw[tunnel] (v4s) -- (v4e)
  node[midway, above, font=\small\sffamily\bfseries, text=clAccent]
    {тунель (UDP\,:51820)};
\coordinate (v4rs) at ($(W.south)+(0,-9.1cm)$);
\coordinate (v4re) at ($(C.south)+(0,-9.1cm)$);
\draw[tunnel] (v4rs) -- (v4re);
%%
\end{tikzpicture}
